\documentclass[10pt,a4paper]{amsart}
\usepackage[margin=19mm]{geometry}
\usepackage{amsmath,amssymb,mathtools}
\usepackage[scr=rsfs,cal=euler]{mathalfa}
\usepackage{xfrac}
\usepackage[unicode,hidelinks,bookmarks=false]{hyperref}
\newcommand{\ie}{i.e.\ }
\newcommand{\cf}{cf.\ }
\newcommand{\resp}{respectively\ }
\newcommand{\Subsection}{\subsection}
\providecommand{\nobreakdash}{\nobreak\hbox{-}}
\setlength{\emergencystretch}{2em}
\multlinegap0pt
\usepackage{bm}
\usepackage{bbm}
\usepackage{dsfont}
\usepackage{xspace}
\usepackage{cancel}
\usepackage{mathtools}
\usepackage{amsthm}
\makeatletter
\@ifundefined{lem}{\newtheorem{lem}{Lemma}}{}
\@ifundefined{prop}{\newtheorem{prop}{Proposition}}{}
\makeatother
\newcommand{\RN}[1]{%
  \textup{\uppercase\expandafter{\romannumeral#1}}%
}
\title[Capillary minimal slicing and scalar curvature rigidity]{Capillary minimal slicing and scalar curvature rigidity}
\author{Dongyeong Ko, Xuan Yao}
\date{Source: arXiv:2602.21071v1 (CC BY 4.0)}
\begin{document}
\maketitle

\noindent\emph{Source and licence.} Capillary minimal slicing and scalar curvature rigidity. Version arXiv:2602.21071v1 (CC BY 4.0), distributed under the Creative Commons Attribution 4.0 International licence, as recorded in the arXiv licence metadata for this exact version and shown on the abstract page https://arxiv.org/abs/2602.21071v1. Full rights notice: licenses/arxiv-2602.21071-CC-BY-4.0.txt.\par\smallskip

\noindent\textit{Editorial scope.} Counted unit: the Prop whose source label is \texttt{prop: winding number} in the article (source0.tex lines 470--475, source label \texttt{prop: winding number}), delivered together with the article's own complete original proof of it (source0.tex lines 476--627, one \texttt{proof} environment of 152 lines). No mathematics is written, repaired or re-derived: statement and proof are the article's own text line for line. Required context retained uncounted: prop: Angle Identities (lines 364--370); prop:local estimates (lines 420--448); lem2.1diag (lines 435--439); threepositivemccomparison (lines 435--439). Every definition, display and label of those lines is retained unchanged. Omitted from this package and not redistributed: 1--363, 371--419, 440--469, 628--1165. Nothing inside a retained excerpt is omitted or altered: every retained body is delivered exactly as the article's lines are written, so no omission receipt of any kind accompanies this package. Citations: the retained excerpts carry no citation command at all, so no bibliography entry of the article is transcribed and no reference is redistributed. Reference adaptation: none was needed and none was made; every reference inside the delivered unit resolves inside it, so no unresolved reference or citation is produced. Printed slips kept literal: no printed word, symbol, hypothesis or number of the counted statement or of its proof was altered. Layout and class: the article's own class is \texttt{amsart} with packages including amsmath, amssymb, mathrsfs, xcolor, graphicx, appendix, comment, ulem, geometry, url; the wrapper is the lane's pinned \texttt{amsart} editorial wrapper at 10pt on A4, which restates the theorem-like environments the retained text needs and adds the article's own macro definitions that the retained text needs (\texttt{\textbackslash RN}), with extra wrapper packages (bm, bbm, dsfont, xspace, cancel, mathtools). The delivered numbering is the unit's own. Assets: no retained span contains a figure environment, a graphics-inclusion command, a PDF-page inclusion, a table or a TikZ picture; no image file is copied into this package, the package declares no asset dependency and no image file is redistributed with it. Licence: version arXiv:2602.21071v1 of this article is distributed under the Creative Commons Attribution 4.0 International licence, as recorded in the arXiv licence metadata for this exact version and shown on the abstract page https://arxiv.org/abs/2602.21071v1. A full rights notice is delivered at licenses/arxiv-2602.21071-CC-BY-4.0.txt. Characters: every retained excerpt is pure ASCII, so the delivered engine, XeLaTeX with its default Latin Modern text font, reports no missing character.
\begin{prop}[Angle Identities]\label{prop: Angle Identities}
Suppose $e$ is any vector in the tangent space $T_p\partial M$, for any $k\in \{1,2,\cdots,n\}$, we have that
\begin{align}
    \Pi_{i=0}^{k-1}\sin\bar{\theta}_i\partial_{e}\bar{\theta}_k=\RN{2}(e,\bar{\nu}_k),
\end{align}
    where we define $\sin\bar{\theta}_0=1$.
\end{prop}

\begin{prop}\label{prop:local estimates}
For any $p\in \partial\Sigma$, suppose that $\{\tau_i\}_{i=1}^n$ is an orthonormal frame of $T_p \partial M$ with respect to $g$. Assume further that $p\in\partial\bar{\Sigma}_t$ for some $t$, and $\{\bar{\tau}_i\}_{i=1}^n$ is an orthonormal frame of $T_p \partial M$, with respect to $g_E$, then
\begin{align}\label{eq: local estimate}
    H_{\partial M}\geq \sum_{i=1}^n\RN{2}(\tau_i,\bar{\tau}_i),
\end{align}
where $H_{\partial M}^{2} g|_{F^{\#}\RN{2}_{+}} = \bar{H}_{\partial M}^{2} g_{E}|_{\RN{2}_{+}}$ when the equality holds. 

\begin{proof}
   Since the inequality is symmetric, we can rotate $\{\bar{\tau}_i\}_{i=1}^n$ appropriately and take the matrix $(a_{ij})_{1 \le i,j \le n} = (\RN{2}(\bar{\tau}_i,\bar{\tau}_j))_{1 \le i,j \le n}$ to be a diagonal matrix. We rotate $\{\tau_i\}_{i=1}^n$ simultaneously with $\{\bar{\tau}_i\}_{i=1}^n$. More specifically there exists an orthogonal matrix $(O_{ij})_{n\times n}\in SO(n)$ such that $\tau_{i}' = O_{ij}\tau_{j}$ and $\bar{\tau}_{i}' = O_{ij} \bar{\tau}_{j}$ and we replace $\tau_{i}$ and $\bar{\tau}_{i}$ by $\tau_{i}'$ and $\bar{\tau}_{i}'$, respectively. Note that $\{\tau_i\}_{i=1}^n$ is preserved to be an orthonormal set in metric $g$ since a rotation is represented by an orthogonal matrix.
 \\
        Let
    \[
    W_i=\frac{a_{ii}}{\bar{H}_{\partial M}}
    \]
    for $i\in\{1,2,\cdots,3\}$ and note that $\sum_{i=1}^nW_i=1$. By the boundary comparison condition, we have
\begin{align}
  \label{threepositivemccomparison} H_{\partial M}\geq& \sum_{i=1}^nW_i\bar{H}_{\partial M} g_E(\tau_i,\bar{\tau}_{i}) \\
  \nonumber &= \sum_{i=1}^na_{ii} g_E(\tau_i,\bar{\tau}_{i})\\
    &=\sum_{i=1}^n\RN{2}(\tau_i,\bar{\tau}_i) \label{lem2.1diag},
    \end{align}
where (\ref{threepositivemccomparison}) gives $H_{\partial M}^{2} g|_{F^{\#}\RN{2}_{+}} = \bar{H}_{\partial M}^{2} g_{E}|_{\RN{2}_{+}}$ (\ref{lem2.1diag}) follows from the coordinate decomposition and the fact that $(a_{ij})$ is a diagonal matrix.
Also note that
\begin{align*}
    \sum_{i=1}^n\RN{2}(O_{ij}\tau_j,O_{ik}\bar{\tau}_k)&= \sum_{i=1}^nO_{ij}O_{ik}\RN{2}(\tau_j,\bar{\tau}_k) \\ &= \delta_{jk} \RN{2}(\tau_j,\bar{\tau}_k) \\ &= \sum_{i=1}^n\RN{2}(\tau_i,\bar{\tau}_i).
\end{align*}
Hence the right hand side sum is an invariant over rotations. The equality condition also follows from the rotation.
\end{proof}

\end{prop}

\begin{prop} \label{prop: winding number}
\begin{align}
    \int_{\partial\Sigma_n}\frac{1}{\Pi_{i=1}^n\sin\bar{\theta_i}}\left(H_{\partial M}-\sum_{i=1}^n\Pi_{j\leq i}\sin\bar{\theta}_j\partial_{\nu_i}\bar{\theta}_i\right)\geq 2k\pi,
    \end{align}
    for some $k\in\mathbb N_+$.
\end{prop}

\begin{proof}
    By Proposition \ref{prop: Angle Identities} and Proposition \ref{prop:local estimates}, we have that
    \begin{align*}
        \int_{\partial\Sigma_n}\frac{1}{\Pi_{i=1}^n\sin\bar{\theta_i}}\left(H_{\partial M}-\sum_{i=1}^n\Pi_{j\leq i}\sin\bar{\theta}_j\partial_{\nu_i}\bar{\theta}_i\right)dl_g\geq \int_{\partial\Sigma_n}\frac{\RN{2}(\tau_1,\bar{\tau}_1)}{\Pi_{i=1}^n\sin\bar{\theta}_i}dl_g
        \end{align*}
        Since $\bar{M}$ is a convex domain in $\mathbb R^n$, we parametrize $\partial \bar{M}$ as 
        \begin{align}
            \Psi(\vec{u})=\Psi(u_1,\cdots, u_{n+1}):=(x(\vec{u}),y(\vec{u}),u_2,u_3,\cdots, u_{n+1}),
        \end{align}
        and $\bar{\tau}_1$ is parallel to $\frac{\partial\Psi}{\partial u_1}$.
        Under this parametrization, we know that $\bar{X}$ is parallel to
        \begin{align*}
           \mathcal{X}_n=\det \left(\begin{matrix}
                \hat{i}_1 & \hat{i}_2 & \hat{i}_3 &\cdots &\hat{i}_{n+2}\\
                x_{u_1} & y_{u_1} & 0 &\cdots & 0\\
                x_{u_2} & y_{u_2} & 1 & \cdots & 0\\
                x_{u_3} &y_{u_3} &0 &\cdots & 0\\
                \vdots  & \vdots & \vdots &\ddots & \vdots\\
                x_{u_{n+1}} & y_{u_{n+1}} & 0 &\cdots & 1 
            \end{matrix}\right)
        \end{align*}
        Similarly, one define $\mathcal{X}_k$ for $k\in\{2,3,\cdots, n\}$, we use this definition to compute the determinant in a neat way.
        \begin{lem}
            \begin{align}\label{normalvecinduction1}
                \mathcal{X}_{n}=\mathcal{X}_{n-1}+c_{n+2}\hat{i}_{n+2}.
            \end{align}
        \end{lem}
        \begin{proof}
            For any $\hat{i}_{m}$ component of $\mathcal{X}_n$ when $m\in\{1,2,\cdots, n\}$, the coefficient is 
            \begin{align*}
                \det\left(\begin{matrix}
                    M_{n-1,m} & \vec{0}\\
                    \vec{v}_n & 1
                \end{matrix}\right),
            \end{align*}
            where $\vec{0}$ is a $(n-1)\times 1$ column whose entries are all zero,
            \[
            \vec{v}_n=(x_{u_{n+1}}, y_{u_{n+1}}, 0,\cdots, 0),
            \]
            and $\det M_{n-1,m}$ is the coefficient of $\hat{i}_m$ in $\mathcal{X}_{n-1}$. The equation \eqref{normalvecinduction1} follows.
        \end{proof}
        Our next step is to compute $c_n$.

       \begin{lem}
           \begin{align}
               c_{n+2}=x_{u_1}y_{u_{n+1}}-x_{u_{n+1}}y_{u_1}.
           \end{align}
       \end{lem}
       \begin{proof}
           Note that 
           \begin{align*}
               c_{n+2}&=(-1)^{n-1}\left((-1)^{n-2}x_{u_{n+1}}\det\left(\begin{matrix}
                   y_{u_1} & 0 & 0 &\cdots & 0\\
                   y_{u_2} & 1 & 0 & \cdots & 0\\
                   y_{u_3} & 0 & 1 & \cdots & 0\\
                   \vdots & \vdots & \vdots & \ddots & \vdots \\
                   y_{u_{n}} & 0  & 0 &\cdots& 1
               \end{matrix}\right)+(-1)^{n-1}y_{u_{n+1}}\det\left(\begin{matrix}
                   x_{u_1} & 0 & 0 &\cdots & 0\\
                   x_{u_2} & 1 & 0 & \cdots & 0\\
                   x_{u_3} & 0 & 1 & \cdots & 0\\
                   \vdots & \vdots & \vdots & \ddots & \vdots \\
                   x_{u_{n}} & 0  & 0 &\cdots& 1
               \end{matrix}\right)\right)\\
               &=x_{u_1}y_{u_{n+1}}-x_{u_{n+1}}y_{u_1}.
           \end{align*}
       \end{proof}
    We then have that
    \begin{align}
        \mathcal{X}_n=&\mathcal{X}_1+\sum_{m=3}^{n+2}\left(x_{u_1}y_{u_{m-1}}-x_{u_{m-1}}y_{u_1}\right)\hat{i}_m\\
        \nonumber=&y_{u_1}\hat{i}_1-x_{u_1}\hat{i}_2+\sum_{m=3}^{n+2}\left(x_{u_1}y_{u_{m-1}}-x_{u_{m-1}}y_{u_1}\right)\hat{i}_m.    \end{align}

    Since $\bar{X}$ is orthogonal to $\bar{\tau}_1$ in the Euclidean metric, we have
    \begin{align*}
        g_E(\nabla_{\gamma'(t)}\bar{X},\bar{\tau}_1)&=\frac{1}{|\mathcal{X}_n|}g_E(\frac{d}{dt}\mathcal{X}_n,\bar{\tau}_1)\\
        &=\frac{\sum_{i=1}^{n+1}x_{u_1}y_{u_1u_i}u_i'-\sum_{i=1}^{n+1}y_{u_1}x_{u_1u_{i}}u_i'}{|\mathcal{X}_n|}\\
        &=\frac{\sum_{i=1}^{n+1}x_{u_1}y_{u_1u_i}u_i'-\sum_{i=1}^{n+1}y_{u_1}x_{u_1u_{i}}u_i'}{\sqrt{y_{u_1}^2+x_{u_1}^2+\sum_{m=3}^{n+2}(x_{u_1}y_{u_{m-1}}-x_{u_{m-1}}y_{u_1})^2}\sqrt{x^2_{u_1}+y^2_{u_1}}}.
    \end{align*}

\begin{lem}
    For $m\in\{1,2\cdots, n\}$, we have that
    \begin{align}\label{eq:angle induction}
        \Pi_{i=1}^m\sin\bar{\theta}_i=\frac{|\mathcal{X}_{n-m}|}{|\mathcal{X}_{n}|}.
    \end{align}
\end{lem}
\begin{proof}

We prove by induction. For $m=1$, we have that
\begin{align*}
    \cos\bar{\theta}_1=&g_E(\bar{X},e_n)\\
                      =&\frac{1}{|\mathcal{X}_n|}g_E(\mathcal{X}_n,e_n)\\
                      =&\frac{c_n}{|\mathcal{X}_n|}.
\end{align*}

Now we have that
\begin{align*}
    \sin\bar{\theta}_1=&\sqrt{1-\cos^2\bar{\theta}_1}\\
                      =&\sqrt{1-\frac{c_n^2}{|\mathcal{X}_n|^2}}\\
                      =&\sqrt{1-\frac{|\mathcal{X}_n|^2-|\mathcal{X}_{n-1}|^2}{|\mathcal{X}_n|^2}}\\
                      =&\frac{|\mathcal{X}_{n-1}|}{|\mathcal{X}_n|},
\end{align*}
where we used \eqref{normalvecinduction1}.

We verified \eqref{eq:angle induction} for $m=1$. Assume that $\eqref{eq:angle induction}$ is true for $1\leq i\leq m$, we consider the case when $i=m+1$.

Since for any $m\in\{1,2,\cdots, n-1\}$, we have that
\begin{align*}
    \bar{\eta}_{m}=\frac{1}{\sin\bar{\theta}_m}\left(\bar{\eta}_{m-1}-\cos\bar{\theta}_{m}e_{n+2-m}\right).
\end{align*}

By definition 
\begin{align*}
    \cos\bar{\theta}_{m+1}=&g_E(\bar{\eta}_{m},e_{n+1-m})\\
                          =&\frac{1}{\sin\bar{\theta}_m}g_E(\bar{\eta}_{m-1},e_{n+1-m})
\end{align*}
Observe that 
\begin{align*}
    \bar{\eta}_{m-1}=\frac{1}{\sin\bar{\theta}_{m-1}}\left(\bar{\eta}_{m-2}-\cos\bar{\theta}_{m-1}e_{n+3-m}\right),
\end{align*}
where the $e_{n+2-m}$ component is perpendicular to $e_{n+1-m}$, and thus we have that
\begin{align*}
    \cos\bar{\theta}_{m+1}=\frac{1}{\sin\bar{\theta}_{m}\sin\bar{\theta}_{m-1}}g_E(\bar{\eta}_{m-2},e_{n+1-m}).
\end{align*}
The above procedure proceeds, and we finally have that
\begin{align*}
    \cos\bar{\theta}_{m+1}=&\frac{g_E(\bar{X},e_{n+1-m})}{\Pi_{i=1}^{m}\sin\bar{\theta}_i}\\
    =&\frac{c_{n+2-m}}{|\mathcal{X}_n|\Pi_{i=1}^m\sin\bar{\theta}_i}\\
    =&\frac{c_{n+2-m}}{|\mathcal{X}_{n-m}|},
\end{align*}
where we used the induction hypothesis in the last step.
We then have that
\begin{align*}
    \sin\bar{\theta}_{m+1}=&\sqrt{1-\frac{c_{n+2-m}^2}{|\mathcal{X}_{n-m}|^2}}\\
    =&\frac{|\mathcal{X}_{n-m-1}|}{|\mathcal{X}_{n-m}|},
\end{align*}
where we used \eqref{normalvecinduction1}.
Now we have 
\[
\Pi_{i=1}^{m+1}\sin\bar{\theta}_i=\frac{|\mathcal{X}_{n-m-1}|}{|\mathcal{X}_n|},
\]
the induction completes.
\end{proof}
Now, we have that
\begin{align*}
    \int_{\partial\Sigma_n}\frac{\RN{2}(\tau_1,\bar{\tau}_1)}{\Pi_{i=1}^{n}\sin\bar{\theta}_i}=&\int_{\partial\Sigma_n}
    \frac{\sum_{i=1}^{n+1}\left(x_{u_1}y_{u_1u_i}u_i'(t)-y_{u_1}x_{u_1u_{i}}u_i'(t)\right)}{|\mathcal{X}_n|\sqrt{x_{u_1}^2+y_{u_1}^2}}\Pi_{i=1}^{n}\sin\bar{\theta}_1dt\\
    =&\int_{\partial\Sigma_n}\frac{\sum_{i=1}^{n+1}\left(x_{u_1}y_{u_1u_i}-y_{u_1}x_{u_1u_i}\right)u_i'(t)}{x_{u_1}^2+y_{u_1}^2}dt\\
    =&\int_{\partial\Sigma_n}\frac{d}{dt}\arctan(\frac{y_{u_1}}{x_{u_1}})dt\\
    =&2k\pi,
\end{align*}
for some $k\geq 1$.
\end{proof}

\end{document}
